\chapter[Differential Theory of the UOD]{Differential Theory of the Universal Optimality Defect}

\section{Introduction}

Chapter 21 established an exact algebraic formula for the Universal
Optimality Defect (UOD). The present chapter develops its differential
theory.

Throughout, let

\[
F=Q\circ L:\mathcal P\rightarrow F(\mathcal P)
\]

be the global logarithmic embedding of Part~III and

\[
\mathcal D(P,Q)=\|F(P)-F(Q)\|^{2}.
\]

Since \(F\) is a global diffeomorphism and the pullback metric is defined
by \(g=F^{*}g_E\), all differential properties of the UOD can be derived
directly from Euclidean calculus.

\section{Differential of the Logarithmic Embedding}

For \(P\in\mathcal P\),

\[
DF_P=Q\circ DL_P.
\]

By Part~III, \(DF_P\) is a linear isomorphism on every tangent space.
Consequently,

\[
g_P(u,v)=\langle DF_Pu,DF_Pv\rangle.
\]

\section{First Variation}

\begin{theorem}[First Variation of the UOD]
\label{ch:thm:22-1}
For every smooth curve \(P(t)\),

\[
\frac{d}{dt}\mathcal D(P(t),Q)=
2\langle DF_{P(t)}P'(t),F(P(t))-F(Q)\rangle.
\]
\end{theorem}

\begin{proof}
Differentiate the Euclidean squared norm and apply the chain rule.
$\blacksquare$
\end{proof}

\section{Characterization of Critical Points}

\begin{theorem}[Unique Critical Point]
\label{ch:thm:22-2}
The only critical point of

\[
P\mapsto\mathcal D(P,Q)
\]

is \(P=Q\).
\end{theorem}

\begin{proof}
If \(P\) is critical then \(d\mathcal D_P(v)=0\) for every tangent
vector \(v\). By Theorem~\ref{ch:thm:22-1},

\[
\langle DF_Pv,F(P)-F(Q)\rangle=0
\]

for every \(v\). Since \(DF_P\) is an isomorphism, its image is the
whole tangent space of \(F(\mathcal P)\). Hence \(F(P)-F(Q)=0\).
The injectivity of \(F\) (Part~III) implies \(P=Q\). $\blacksquare$
\end{proof}

\section{Gradient}

\begin{corollary}[Riemannian Gradient]
\label{ch:cor:22-3}
The Riemannian gradient is uniquely determined by

\[
g_P(\nabla_g\mathcal D,v)=2\langle DF_Pv,F(P)-F(Q)\rangle.
\]

Equivalently,

\[
\nabla_g\mathcal D = 2\,DF_P^{-1}(F(P)-F(Q)).
\]
\end{corollary}

\begin{proof}
Using the pullback metric \(g_P(u,v)=\langle DF_Pu,DF_Pv\rangle\), the
defining identity of the Riemannian gradient gives

\[
DF_P(\nabla_g\mathcal D)=2(F(P)-F(Q)).
\]

Invertibility of \(DF_P\) completes the proof. $\blacksquare$
\end{proof}

\section{Hessian at the Reference Posterior}

\begin{theorem}[Hessian at the Optimum]
\label{ch:thm:22-4}
At the reference point \(P=Q\),

\[
\operatorname{Hess}_g(\mathcal D)_P = 2I.
\]
\end{theorem}

\begin{proof}
Near \(P\),

\[
F(P+\delta) = F(P)+DF_P\delta+o(\|\delta\|).
\]

Therefore,

\[
\mathcal D(P+\delta,P) = \|DF_P\delta\|^{2}+o(\|\delta\|^{2}).
\]

Using the definition of the pullback metric,
\(\|DF_P\delta\|^{2}=g_P(\delta,\delta)\), so

\[
\mathcal D(P+\delta,P) = g_P(\delta,\delta)+o(\|\delta\|^{2}).
\]

Hence the quadratic term is represented by the identity bilinear form,
whose Hessian is \(2I\). $\blacksquare$
\end{proof}

\section{Strict Convexity Near the Optimum}

\begin{corollary}[Local Convexity]
\label{ch:cor:22-5}
The reference posterior is a strict local minimizer of the UOD.
\end{corollary}

\begin{proof}
The Hessian at \(P\) equals \(2I\), which is positive definite.
$\blacksquare$
\end{proof}

\section{Local Quadratic Model}

Combining the previous results,

\[
\boxed{
\mathcal D(P+\delta,P)
=
\|\delta\|_g^2
+
o(\|\delta\|_g^2).
}
\]

Thus the UOD is locally identical to the squared Riemannian distance
induced by the pullback metric.

\section{Consequences}

The differential theory establishes:
\begin{itemize}
\item smoothness of the UOD;
\item uniqueness of its critical point;
\item explicit gradient;
\item exact Hessian at the optimum;
\item strict local convexity;
\item local quadratic approximation.
\end{itemize}

These properties are intrinsic consequences of the logarithmic embedding
and require no additional probabilistic assumptions.

\section{Outlook}

The next chapter compares this local quadratic structure with the
classical Fisher information metric, identifying the precise
relationship between the WIS geometry and standard information geometry.
